\section{Matriz de trazabilidad de escenarios}

\begin{longtable}{p{1.8cm}p{3.2cm}p{4.4cm}p{2.7cm}}
\caption{Trazabilidad requisito--actividad--evidencia}\label{tab:trace}\\
\toprule
Escenario & Actividad & Evidencia principal & Condición de origen \\
\midrule
\endfirsthead
\toprule
Escenario & Actividad & Evidencia principal & Condición de origen \\
\midrule
\endhead
1 & Análisis del caso & \texttt{docs/}\allowbreak\texttt{analisis\_inicial.md} & Literal \\
2 & Clasificación ISO 25022 & CSV clasificado y resumen JSON & Literal \\
3 & Niveles de calidad y SQuaRE & Marco teórico y Figura \ref{fig:square} & Literal \\
4 & Atributos UTC & Tabla \ref{tab:utc} & Reconstruido \\
5 & Mapeo y prioridad & Matriz de tareas críticas & Reconstruido \\
6 & Diseño de métricas & Catálogo, fórmulas y Tabla \ref{tab:catalog} & Reconstruido \\
7 & Obtención de datos & Eventos, encuestas y manifiesto & Reconstruido \\
8 & Automatización & Paquetes de \texttt{analytics} & Reconstruido \\
9 & KPI y tablero & Tabla \ref{tab:results} y React gerencial & Reconstruido \\
10 & Interpretación & Hipótesis de causa y prioridades & Reconstruido \\
11 & Comunicación ejecutiva & Síntesis de decisiones & Reconstruido \\
12 & Mejora continua & Plan PDCA, Tabla \ref{tab:pdca} & Reconstruido \\
Reto & Telemedicina 2.0 & Trazabilidad y casos de contexto & Reconstruido \\
\bottomrule
\end{longtable}

\section{Ubicación de artefactos}

\begin{table}[H]
\centering
\caption{Estructura de evidencia reproducible}
\begin{tabularx}{\textwidth}{p{4.0cm}Y}
\toprule
Ruta & Responsabilidad \\
\midrule
\texttt{docs/Documento}\allowbreak\texttt{ Padre} & Fuente académica principal y dataset recibido. \\
\texttt{data/original} & Copia preservada del dataset de 3.000 incidentes. \\
\texttt{data/simulated} & Eventos, encuestas, manifiestos y ejecuciones regenerables. \\
\texttt{data/processed} & Clasificación y métricas derivadas. \\
\texttt{analytics} & Clasificación, simulación, medición y exportación en Python. \\
\texttt{apps/mobile-patient} & Portal Flutter del paciente. \\
\texttt{apps/web-his} & HIS operativo React y módulo gerencial de calidad. \\
\texttt{services} & Gateway, dominios clínicos y API de medición. \\
\texttt{evidencias/}\allowbreak\texttt{resultados} & Resúmenes seleccionados para entrega. \\
\texttt{informe/entrega} & PDF académico final. \\
\bottomrule
\end{tabularx}
\end{table}

\section{Controles de reproducción}

\begin{enumerate}
  \item Verificar el hash SHA-256 del dataset original contra el archivo del Documento Padre.
  \item Instalar las dependencias de Python y ejecutar el pipeline con semilla 25022.
  \item Confirmar 3.000 clasificaciones, 6.000 eventos, 2.000 encuestas y 10 métricas.
  \item Repetir la simulación y comparar manifiestos y resultados.
  \item Levantar Docker Compose desde Ubuntu WSL y comprobar los estados de salud.
  \item Probar React y Flutter contra el mismo Gateway local.
  \item Compilar este documento con XeLaTeX mediante \texttt{latexmk}.
  \item Renderizar todas las páginas y revisar portada, tablas, figuras, referencias y desbordamientos.
\end{enumerate}

\section{Criterio de procedencia}

Para evitar ambigüedad, todo resultado debe identificarse con una de estas etiquetas:

\begin{description}[style=nextline,leftmargin=2.5cm,labelwidth=2.2cm]
  \item[Original:] contenido recibido en el Documento Padre o en su dataset asociado.
  \item[Simulado:] observación generada artificialmente con semilla y reglas declaradas.
  \item[Derivado:] clasificación, agregado, fórmula, gráfica o interpretación calculada a partir de una fuente.
  \item[Propuesto:] métrica, meta, hipótesis o acción definida por el equipo para completar el taller.
\end{description}
